\section{Dynamic Preference Learning}
\label{sec:learning}

The proposed updater aggregates signed, timestamped evidence by dimension and
scope. Explicit dashboard changes and direct statements have greater authority
than repeated implicit corrections; one ambiguous event cannot create a durable
high-confidence preference. A newer direct durable statement can supersede an
unlocked inferred value, while weaker contradictory evidence maintains an
explicit ambiguous state and causes selection to abstain. Locked values do not
change automatically, but the conflict remains inspectable.

Recency decay is applied to inferential evidence, while user-locked preferences do
not decay. Temporary, task-local requests affect the current generation but are
not promoted to durable profile values without an explicit persistence cue. Drift
is evaluated through stable change points, temporary overrides, and alternating
task-dependent preferences.

The reference interface also accepts negative applicability evidence: ``this
preference was not relevant here'' adds a scoped exception while retaining the
underlying preference assertion. This creates direct supervision for the second
error type without conflating it with preference truth.

\evidencegap{Learning curves, activation precision/recall, calibration, and
adaptation-delay estimates require the frozen longitudinal experiments.}
